\documentclass[11pt]{article}

\usepackage[T1]{fontenc}
\usepackage{amsmath,amssymb,amsthm}
\usepackage{hyperref}

\title{Model-Relative Bias of Omega Under Equivalent Measurement Models}
\author{}
\date{\today}

\newcommand{\E}{\mathbb{E}}
\newcommand{\tr}{\operatorname{tr}}
\newcommand{\Var}{\operatorname{Var}}
\newcommand{\Cov}{\operatorname{Cov}}
\newcommand{\vech}{\operatorname{vech}}
\newcommand{\argmin}{\operatorname*{arg\,min}}
\newcommand{\argmax}{\operatorname*{arg\,max}}
\newcommand{\plim}{\operatorname*{plim}}
\newcommand{\T}{^{\mathsf T}}

\theoremstyle{plain}
\newtheorem{theorem}{Theorem}
\newtheorem{lemma}{Lemma}
\newtheorem{corollary}{Corollary}
\theoremstyle{remark}
\newtheorem*{remark}{Remark}

\begin{document}
\maketitle

\section*{Purpose}

Bell, Chalmers, and Flora \cite{bell2024} define population reliability as the
omega parameter attached to the model that generated the covariance matrix.
They then call an estimated omega biased when its expectation differs from that
model-relative target. That is internally coherent.

The problem is that the target is not, in general, a functional of the observed
item distribution. Equivalent latent decompositions can imply the same
covariance matrix while assigning different amounts of composite variance to the
named target factor. In that case the same omega estimate has different
population bias under different, observationally indistinguishable, ``true''
models. This note gives a simple theorem and a numerical example where the bias
sign changes without changing the observable covariance matrix.

\section{Model-Relative Omega}

Let \(y\in\mathbb{R}^p\) be the item vector and let \(X=w\T y\) be the composite,
usually with \(w=\mathbf{1}\). Write \(\Sigma=\Var(y)\). If a target factor \(f\)
with \(\Var(f)=1\) contributes the item loading vector \(a\), then the omega
target for the composite is
\begin{equation}\label{eq:omega}
  \omega(a,\Sigma;w)
  =
  \frac{(w\T a)^2}{w\T\Sigma w}.
\end{equation}
For a one-factor model, \(a\) is the loading vector. For a bifactor model, \(a\)
is the general-factor loading vector. For a higher-order model, \(a\) is the
vector of indirect effects from the higher-order factor to the items. The
Bell et al. bias calculation is therefore
\[
  \E(\hat\omega)-\omega(a_{\mathrm{true}},\Sigma;w),
\]
where \(a_{\mathrm{true}}\) is supplied by the chosen population model.

The word ``chosen'' matters. The covariance matrix alone need not determine
\(a_{\mathrm{true}}\).

\section{A Non-Identification Result}

\begin{theorem}[Equivalent models with arbitrary composite omega]
\label{thm:anyomega}
Let \(\Sigma\) be positive definite and let \(w\ne 0\). For every
\(r\in[0,1]\), define
\begin{equation}\label{eq:ar}
  a_r
  =
  \sqrt{r}\,
  \frac{\Sigma w}{\sqrt{w\T\Sigma w}},
  \qquad
  \Psi_r
  =
  \Sigma-a_r a_r\T .
\end{equation}
Then \(\Psi_r\succeq0\), the one-factor representation
\[
  y=a_r f+u,\qquad \Var(f)=1,\qquad \Cov(f,u)=0,\qquad \Var(u)=\Psi_r
\]
implies the same covariance matrix \(\Sigma\), and
\[
  \omega(a_r,\Sigma;w)=r .
\]
If all entries of \(\Sigma w\) are nonzero, the target factor loads on every
item for every \(r>0\).
\end{theorem}

\begin{proof}
The covariance equality is immediate:
\[
  \Var(a_r f+u)=a_r a_r\T+\Psi_r=\Sigma .
\]
For any vector \(x\), Cauchy's inequality in the inner product
\(\langle x,z\rangle_\Sigma=x\T\Sigma z\) gives
\[
  (x\T\Sigma w)^2
  \le
  (x\T\Sigma x)(w\T\Sigma w).
\]
Therefore
\[
  x\T\Psi_r x
  =
  x\T\Sigma x
  -
  r\frac{(x\T\Sigma w)^2}{w\T\Sigma w}
  \ge
  (1-r)x\T\Sigma x
  \ge 0 ,
\]
so \(\Psi_r\succeq0\). Finally,
\[
  (w\T a_r)^2
  =
  r\,w\T\Sigma w ,
\]
which proves \(\omega(a_r,\Sigma;w)=r\).
\end{proof}

\begin{corollary}[Bias sign is not identified]
\label{cor:bias}
Let an omega estimator converge under the observed covariance \(\Sigma\) to
some value \(T(\Sigma)\in(0,1)\). In the equivalent model family from
Theorem~\ref{thm:anyomega}, its asymptotic bias is
\[
  T(\Sigma)-r .
\]
Thus, without changing the observed covariance matrix, one can choose an
equivalent population model under which the estimator is positively biased,
unbiased, or negatively biased.
\end{corollary}

The corollary is the clean version of the objection. Bias is identifiable only
after fixing a target factor by theory or by a restrictive model class that makes
the target unique. It is not identified by the covariance matrix itself.

\section{Numerical Example}

Take \(p=8\), \(w=\mathbf{1}\), and the compound-symmetry correlation matrix
\[
  \Sigma_\rho=(1-\rho)I+\rho\mathbf{1}\mathbf{1}\T,
  \qquad
  \rho=.30 .
\]
The independent-error one-factor decomposition has
\[
  a_0=\sqrt{\rho}\,\mathbf{1},
  \qquad
  \Psi_0=(1-\rho)I .
\]
Its omega is
\[
  \omega_0
  =
  \frac{p^2\rho}{p+p(p-1)\rho}
  =
  0.774194 .
\]
A correctly fitted one-factor omega also converges to \(T(\Sigma_\rho)=0.774194\).

But the same \(\Sigma_\rho\) has the following equivalent one-factor
representations with correlated residuals:
\[
\begin{array}{c|c|c|c|c}
  r & a_r\text{ entry} & \Psi_r\text{ diagonal} &
  \Psi_r\text{ off diagonal} & T(\Sigma_\rho)-r\\
  \hline
  .600000 & .482183 & .767500 & .067500 & .174194\\
  .774194 & .547723 & .700000 & .000000 & .000000\\
  .850000 & .573912 & .670625 & -.029375 & -.075806
\end{array}
\]
All three rows generate exactly the same item covariance matrix. The first says
the same fitted omega is positively biased by about .17. The second says it is
unbiased. The third says it is negatively biased by about .08. The only thing
that changed is the latent story used to split \(\Sigma_\rho\) into target-factor
variance and residual covariance.

\section{What This Does and Does Not Refute}

The theorem does not show that Bell et al. made a computational mistake. Once
their population model and target factor are stipulated, their Monte Carlo bias
quantity is a legitimate model-relative quantity.

It does show that the broader language ``omega is biased for reliability under
misspecification'' needs a qualification. If reliability means the proportion of
composite variance attributable to a substantively named factor, then that
factor is part of the estimand. Equivalent measurement models that change the
factor contribution vector \(a\) change the estimand while leaving the data
distribution fixed. No sample size or fit index can recover which bias statement
is the correct one from \(\Sigma\) alone.

There is also an important non-objection. Some equivalent representations
preserve the target contribution vector. For example, the Schmid-Leiman
transformation relates restricted bifactor and higher-order models
\cite{schmid1957,yung1999}. When the bifactor general-factor loadings equal the
higher-order indirect effects, the numerator in \eqref{eq:omega} is the same,
so omega is invariant. Equivalence is damaging only when it is not
target-preserving.

\section{Implication for Simulation Design}

A simulation can still answer a conditional question:
\[
  \text{If this latent model is the true target story, what does this fitted
  omega estimate?}
\]
It cannot, by itself, answer an unconditional question:
\[
  \text{How biased is omega for the reliability encoded in the observed
  covariance matrix?}
\]
There is no such covariance-only reliability target unless additional
identifying restrictions are imposed.

For future reliability simulations, the defensible formulation is therefore:
\begin{quote}
  Bias is relative to a declared target model. Before interpreting it as a
  failure of an estimator, show that the target is invariant across the relevant
  equivalent models, or explicitly treat the result as a sensitivity analysis
  over plausible target decompositions.
\end{quote}

Under that standard, the user's objection is right in the strong sense: there
are equivalent ``true'' models under which the same fitted omega has smaller,
larger, zero, or opposite-signed bias. The objection does not make Bell et al.'s
conditional Monte Carlo invalid, but it blocks a covariance-only or
model-free reading of their bias claims.

\begin{thebibliography}{9}

\bibitem{bell2024}
Bell, S. M., Chalmers, R. P., \& Flora, D. B. (2024).
The impact of measurement model misspecification on coefficient omega estimates
of composite reliability.
\emph{Educational and Psychological Measurement}, 84(1), 5--39.
\href{https://doi.org/10.1177/00131644231155804}{doi:10.1177/00131644231155804}.

\bibitem{schmid1957}
Schmid, J., \& Leiman, J. M. (1957).
The development of hierarchical factor solutions.
\emph{Psychometrika}, 22(1), 83--90.
\href{https://doi.org/10.1007/BF02289209}{doi:10.1007/BF02289209}.

\bibitem{yung1999}
Yung, Y. F., Thissen, D., \& McLeod, L. D. (1999).
On the relationship between the higher-order factor model and the hierarchical
factor model.
\emph{Psychometrika}, 64(2), 113--128.
\href{https://doi.org/10.1007/BF02294531}{doi:10.1007/BF02294531}.

\end{thebibliography}

\end{document}
